\documentclass[a4paper,12pt]{article}
\usepackage[utf8]{inputenc}
%\usepackage{ctex}
\usepackage{graphicx}
\usepackage{hyperref}
\usepackage{geometry}
\usepackage{listings}
\usepackage{xcolor}
\usepackage{float}

\geometry{a4paper, margin=1in}

\title{\textbf{RAFTcorr User Manual v1.2}}
\author{Zixiang (Zach) Tong}
\date{\today}

\begin{document}

\maketitle
\tableofcontents
\newpage

\section{Introduction}
RAFTcorr is a 2D Digital Image Correlation (DIC) software powered by the RAFT (Recurrent All-Pairs Field Transforms) optical flow network. It provides a robust, CUDA-accelerated workflow for calculating full-field displacement and strain from image sequences.

\section{Installation}
\subsection{Prerequisites}
\begin{itemize}
    \item Operating System: Windows 10/11
    \item Python 3.8+
    \item \textbf{GPU Requirement}: An NVIDIA GPU is mandatory. The RAFT architecture relies heavily on specific CUDA kernels (Correlation Layer) that are not available on CPU.
    \item \textbf{CUDA Support}: CUDA 11.8+ is recommended for optimal compatibility with modern PyTorch.
    \item \textbf{Numba}: Required for optimized strain calculation (installed automatically via \texttt{pip install}).
\end{itemize}

\subsection{Setup}
\begin{enumerate}
    \item Clone the repository.
    \item Install dependencies: \texttt{pip install -r requirements.txt}
    \item Verify Installation: \texttt{python verify$\_$installation.py}
    \item Run the application: \texttt{python main$\_$GUI.py}
\end{enumerate}

\section{Interface Overview}
The graphical user interface is divided into two main panels:
\begin{enumerate}
    \item \textbf{Control Panel} (Left): For setting parameters, choosing inputs, and controlling the process.
    \item \textbf{Preview Panel} (Right): For ROI selection, visualizing results, and plotting data.
\end{enumerate}

%%%%%%%%%%%%%%%%%%%%%%%%%%%
\begin{figure}[htbp]
    \centering
    \includegraphics[width=1\textwidth]{Figures_manual/fig_manual_overall.bmp}
    \caption{Overview of the RAFTcorr Interface}
    \label{fig:gui_overview}
\end{figure}
%%%%%%%%%%%%%%%%%%%%%%%%%%%

\section{Workflow Guide}

\subsection{Step 1: Image Sequence Selection}
In the \textbf{Path Settings} section of the Control Panel: Click \textbf{Browse} next to "Input Directory" to select the folder containing your image sequence (.tif, .png, .jpg).

\subsection{Step 2: Model Selection}
Select the appropriate RAFT model from the \textbf{Model} dropdown.
\begin{itemize}
    \item \textbf{RAFT-Large}: A downsampling-based model that can usually process larger images in a single pass.  
    Its coarse feature representation greatly reduces memory usage and computation time, making it faster and more scalable for large FOVs.  
    However, the downsampling also limits its spatial precision, resulting in lower accuracy compared with the fine model.

    \item \textbf{RAFT-Fine}: A full-resolution model without any downsampling.  
    It preserves pixel-level details and achieves higher displacement accuracy, especially for small or localized deformations.  
    Due to the significantly higher memory cost, large images must be tiled during inference, leading to slower runtime.
\end{itemize}


\subsection{Step 3: ROI Selection}
Before running, you must define the Region of Interest (ROI) on the reference image in the Preview Panel.
\begin{itemize}
    \item \textbf{Rectangle Tool}: Click and drag to draw a box.
    \item \textbf{Polygon Tool}: Click points to define an irregular shape, double-click to close.
    \item \textbf{Cut (Crop)}: Use the Cut tool to create holes or exclude defects from the mask (e.g., removing a bolt hole).
    \item Click \textbf{Confirm ROI} to finalize the mask.
\end{itemize}

%%%%%%%%%%%%%%%%%%%%%%%%%%%
\begin{figure}[htbp]
    \centering
    \includegraphics[width=1\textwidth]{Figures_manual/fig_manual_ROI.bmp}
    \caption{ROI selection.}
    \label{fig:ROI}
\end{figure}
%%%%%%%%%%%%%%%%%%%%%%%%%%%

\subsection{Step 4: Processing Parameters}
\begin{itemize}
    \item \textbf{Mode}:
    \begin{itemize}
        \item \textbf{Accumulative}: Every frame is correlated against the first (reference) frame. Best for small-to-moderate deformations.
        \item \textbf{Incremental} (New in v1.1): For large deformations, specify key frames where the reference is updated. Displacements are accumulated in original coordinates. Enter comma-separated frame numbers (e.g., ``1, 50, 100'').
    \end{itemize}
    \item \textbf{Smoothing}: Enable Gaussian smoothing and set \textbf{Sigma} (default 2.0) to filter high-frequency noise from the displacement field.
    \item \textbf{Advanced - Tiling}: These settings handle processing of large images that exceed GPU memory.
    \begin{itemize}
        \item \textbf{Safety Factor} (0.2 - 1.0): Controls the memory budget. Lower values (e.g. 0.4) are more conservative and force the image to be split into more tiles. Higher values (0.8+) attempt to process larger chunks. Defaults to 0.55.
        \item \textbf{Context Padding}: Extra pixels (default 32px) fetched around each tile to provide context for the neural network.
        \item \textbf{Tile Overlap}: The pixel width (default 32px) of the blending region between adjacent tiles. This ensures smooth transitions and eliminates grid artifacts.
        \item \textit{Note}: The status text box below the ROI preview will clearly indicate if "Tiling" or "Single-Shot" mode is active based on these parameters and your ROI size.
    \end{itemize}
\end{itemize}

\subsection{Step 5: Run Analysis}
Click the \textbf{Run} button. The progress bar will indicate the status.
\begin{itemize}
    \item The Preview Panel will update in real-time with the displacement field.
    \item Logs can be viewed in the terminal or the Log section.
\end{itemize}

\subsection{Step 6: Displacement Visualization Settings}
\textit{After} the analysis finishes (or during a pause), you can adjust how the displacement results are displayed using the \textbf{Visualization Settings} panel:

%%%%%%%%%%%%%%%%%%%%%%%%%%%
\begin{figure}[htbp]
    \centering
    \includegraphics[width=1\textwidth]{Figures_manual/fig_manual_disp_results.bmp}
    \caption{Displacement visualization.}
    \label{fig:disp_results}
\end{figure}
%%%%%%%%%%%%%%%%%%%%%%%%%%%

\begin{itemize}
    \item \textbf{Range Control}:
    \begin{itemize}
        \item \textbf{Use Fixed Range}: Check to manually lock the Min/Max values for U and V.
        \item \textbf{Range from Frame}: Enter a frame number and click "Apply" to auto-set the range based on that specific time step.
    \end{itemize}
    \item \textbf{Appearance}:
    \begin{itemize}
        \item \textbf{Colormap}: Select a color scheme (e.g., turbo, jet).
        \item \textbf{Transparency}: Adjust overlay opacity (0-1).
        \item \textbf{Background Image Mode} (New in v1.1): Toggle between:
        \begin{itemize}
            \item \textbf{Reference Image}: Data overlaid on the static reference frame.
            \item \textbf{Deformed Image}: Data overlaid on the corresponding deformed frame, showing where material points actually move. Useful for validating tracking accuracy.
        \end{itemize}
    \end{itemize}
    \item \textbf{Physical Units}:
    \begin{itemize}
        \item Enable to scale results from pixels to physical units (e.g., mm).
        \item \textbf{Ratio}: The scaling factor (unit/px).
        \item \textbf{dt(s)}: Time step between frames (for velocity/strain rate calculations).
    \end{itemize}
    \item Click \textbf{Update} to apply changes to the current view.
\end{itemize}

\section{Post-Processing Analysis}
Switch to the \textbf{Post-Processing} tab to calculate strain and perform detailed analysis.

\subsection{Strain Calculation}
To calculate strain fields from the displacement data settings:



\begin{enumerate}
    \item \textbf{Method}:
    \begin{itemize}
        \item \textit{Green-Lagrange}: best for large deformations (includes non-linear geometric terms).
        \item \textit{Engineering}: linear approximation, suitable only for small strains (<1-2\%).
    \end{itemize}
    \item \textbf{VSG Size (px)}: Virtual Strain Gauge size. Defines the local window size for strain fitting. Larger values (e.g. 51px) produce smoother but less localized strain fields. Smaller values (e.g. 31px) preserve gradients but are noisier.
    \item \textbf{Step (px)}: Fixed to 1 for now. Strain is calculated at every pixel for maximum resolution.
    \item \textbf{Poly Order}: The order of the polynomial shape function used to fit the local displacement field (1 = Linear, 2 = Quadratic).
    \item \textbf{Weighting}: 
    \begin{itemize}
        \item \textit{Uniform}: All pixels in the VSG weigh equally.
        \item \textit{Gaussian}: Center pixels have higher weight (recommended regarding smoothness).
    \end{itemize}
    \item \textbf{Components}: Check the tensors to calculate ($e_{xx}$, $e_{yy}$, $e_{xy}$, etc.).
    \item Click \textbf{Calculate Strain}.
\end{enumerate}

\subsection{Strain Visualization}
Once calculated, the "Visualization Controls" section appears:

%%%%%%%%%%%%%%%%%%%%%%%%%%%
\begin{figure}[htbp]
    \centering
    \includegraphics[width=1\textwidth]{Figures_manual/fig_manual_strain.bmp}
    \caption{Strain visualization.}
    \label{fig:strain}
\end{figure}
%%%%%%%%%%%%%%%%%%%%%%%%%%%


\begin{itemize}
    \item \textbf{Display Component}: Choose which strain tensor to view (e.g., $e_{xx}$).
    \item \textbf{Physical Units}: If checked, axes will display physical coordinates based on the Ratio set in Displacement Visualization.
    \item \textbf{Fixed Color Range}: Manually clamp the strain heatmap range (Min/Max) to highlight yield zones or compare samples.
\end{itemize}

\subsection{Probe Analysis}
Use probes to extract quantitative data from the field.
\begin{itemize}
    \item \textbf{Point}: Click multiple spots. Result: Plots the Time vs. Magnitude curve.
    \item \textbf{Line}: Draw multiple lines. Result: Plots the Time vs. Average/Max/Min value along these lines.
    \item \textbf{Area}: Draw multiple regions. Result: Plots the Time vs. Average/Max/Min value within these regions.
\end{itemize}


%%%%%%%%%%%%%%%%%%%%%%%%%%%
\begin{figure}[htbp]
    \centering
    \includegraphics[width=1\textwidth]{Figures_manual/fig_manual_probe_line.bmp}
    \caption{Overview of the RAFTcorr Interface}
    \label{fig:gui_overview}
\end{figure}
%%%%%%%%%%%%%%%%%%%%%%%%%%%

\section{Data Export}
When exporting results via \textbf{Export Scientific Data} (saved as \texttt{.mat} or \texttt{.npz}), the data is organized as flattened variables in the root workspace. 

\subsection{Variable Structure (MATLAB/NumPy)}
Assuming an image sequence of $N$ frames and an ROI size of $H \times W$:

\begin{table}[htbp]
\centering
\begin{tabular}{|l|l|l|}
\hline
\textbf{Variable Name} & \textbf{Dimensions} & \textbf{Description} \\ \hline
\texttt{U}, \texttt{V} & $N \times H \times W$ & Horizontal/Vertical displacement fields. \\ \hline
\texttt{exx}, \texttt{eyy}, \texttt{exy...} & $N \times H \times W$ & Strain tensor components (if calculated). \\ \hline
\texttt{e\_von\_mises} & $N \times H \times W$ & Von Mises equivalent strain. \\ \hline
\texttt{X\_ref}, \texttt{Y\_ref} & $H \times W$ & Pixel coordinates of the Reference ROI (Meshgrid). \\ \hline
\texttt{ROI\_mask} & $H \times W$ & Binary mask defining the Region of Interest. \\ \hline
\texttt{image\_files} & $N \times 1$ & List of filenames for the processed sequence. \\ \hline
\end{tabular}
\caption{Main Field Variables}
\end{table}

\subsection{Metadata Variables}
Scalar parameters recording the analysis settings are also saved:
\begin{itemize}
    \item \texttt{physical\_ratio}, \texttt{physical\_unit}: Unit conversion settings (e.g. 1.0, 'mm').
    \item \texttt{use\_physical\_units}: Boolean flag (1/0).
    \item \texttt{strain\_method}, \texttt{vsg\_size}, \texttt{strain\_step}: Strain calculation parameters.
    \item \texttt{time\_step}: Time interval between frames ($dt$).
\end{itemize}

\subsection{Batch Image Export (New in v1.1)}
Export visualization images for all frames to create figures for publications or animations:
\begin{enumerate}
    \item In the Post-Processing panel, click \textbf{Export Images}.
    \item Select components to export (U, V, $\varepsilon_{xx}$, Velocity, Magnitude, etc.).
    \item Configure options:
    \begin{itemize}
        \item \textbf{Frame Range}: All frames or custom range.
        \item \textbf{Format}: PNG (raster), SVG or PDF (vector).
        \item \textbf{Display Mode}: Reference or Deformed view.
        \item \textbf{Colormap, Alpha, Color Range}: Visual customization.
        \item \textbf{Include Colorbar/Title}: Toggle annotations.
    \end{itemize}
    \item Click \textbf{Export}. Images are saved to a timestamped folder with a settings JSON file for reproducibility.
\end{enumerate}

\subsection{Velocity Visualization (New in v1.2)}
Overlay velocity arrows or streamlines on exported images:
\begin{itemize}
    \item \textbf{Quiver (Arrows)}: Direction-only arrows showing flow direction. All arrows have uniform length; color indicates magnitude.
    \item \textbf{Streamlines}: Continuous flow lines following velocity field.
    \item \textbf{Export Consistency}: Exported images match GUI preview exactly---same arrow density, size, and line width (WYSIWYG).
\end{itemize}

Parameters:
\begin{itemize}
    \item \textbf{Spacing}: Distance between arrows (in pixels). Smaller values = denser arrows.
    \item \textbf{Scale}: Controls arrow size. Larger values = larger arrows.
    \item \textbf{Color}: Arrow/streamline color (white, black, or custom).
    \item \textbf{Line Width}: Thickness of arrows/streamlines.
\end{itemize}

\section{Troubleshooting}
\begin{itemize}
    \item \textbf{OOM (Out of Memory)}: If processing crashes on large images, reduce the \textbf{Safety Factor} (e.g. to 0.4) in the Advanced Tiling section. This forces smaller tiles.
    \item \textbf{Feedback}: We welcome user feedback! Please contact the developer or verify the latest updates on GitHub if you encounter issues.
\end{itemize}

\end{document}
